\documentclass[11pt]{article}
\usepackage[utf8]{inputenc}
\usepackage{amsmath}
\usepackage{amsfonts}
\usepackage{amssymb}
\usepackage{geometry}
\usepackage{booktabs}

\geometry{margin=0.85in}

\title{Absolute Edge-Meter Diagnostics for Synthetic PRISM-FDR Fits}
\author{Jan Balewski, NERSC}
\date{\today}

\begin{document}

\maketitle

\begin{abstract}
\texttt{edgeMaterAbs3c.py} evaluates final aggregate PRISM-FDR connectivity
fits on synthetic data where the true directed graph is known. The evaluated
matrix is \texttt{A\_prune}: the final matrix after bagging, FDR thresholding,
stability selection, source-neuron type classification, and Dale-sign pruning.
The truth edge mask is \texttt{E\_true}; \texttt{A\_true} supplies the true
sign of each true edge. All metrics are computed on directed off-diagonal
pairs only. The default plots \texttt{-p abc} summarize recovery quality,
raw error counts, and selected confusion matrices.
\end{abstract}

\tableofcontents
\newpage

% ===============================================================
\section{Inputs and Edge Labels}
\label{sec:inputs}
% ===============================================================

The script reads one or more aggregate fit files from \texttt{prismFit/},
using a common \texttt{--fitNameTrunk} and a list of \texttt{--fitTags}. Each
fit must point through its provenance metadata to the same synthetic truth
file in \texttt{truthDale/}. The script aborts if different fit tags refer to
different truth graphs.

The matrix convention is source-column:
\[
  A_{ij} = \hbox{effect of source neuron } j
  \hbox{ on target neuron } i .
\]
The candidate graph consists of all off-diagonal directed pairs:
\[
  (i,j), \quad i\ne j .
\]
Diagonal dynamics are useful for the fitted model, but they are not counted
as graph-recovery edges.

For each off-diagonal candidate, the true label is
\[
\mathrm{label}^{\mathrm{true}}_{ij} =
\begin{cases}
+ & E^{\mathrm{true}}_{ij}\ne0 \hbox{ and } A^{\mathrm{true}}_{ij}>0,\\
- & E^{\mathrm{true}}_{ij}\ne0 \hbox{ and } A^{\mathrm{true}}_{ij}<0,\\
0 & E^{\mathrm{true}}_{ij}=0 .
\end{cases}
\]
The recovered label comes from the final \texttt{A\_prune}:
\[
\mathrm{label}^{\mathrm{reco}}_{ij} =
\begin{cases}
+ & A^{\mathrm{prune}}_{ij} > \varepsilon,\\
- & A^{\mathrm{prune}}_{ij} < -\varepsilon,\\
0 & |A^{\mathrm{prune}}_{ij}|\le\varepsilon .
\end{cases}
\]
The default tolerance is \(\varepsilon=10^{-4}\). This keeps the test
effectively binary while avoiding accidental counting of tiny numerical
residue as recovered edges.

% ===============================================================
\section{Plot Group \texttt{-p a}: Metric Curves}
\label{sec:plot-a}
% ===============================================================

Plot group \texttt{a} is the main overview canvas. It contains six panels.
The x-axis is the data duration in minutes when the input fits correspond to
a scan over recording length.

\subsection{Presence recovery}

The presence panel collapses signed labels to a binary edge-present decision.
It shows:
\[
\mathrm{precision}=\frac{\mathrm{TP}}{\mathrm{TP}+\mathrm{FP}},\qquad
\mathrm{recall}=\frac{\mathrm{TP}}{\mathrm{TP}+\mathrm{FN}},
\]
\[
\mathrm{F1}=
\frac{2\,\mathrm{precision}\,\mathrm{recall}}
{\mathrm{precision}+\mathrm{recall}},
\]
and the Matthews correlation coefficient
\[
\mathrm{MCC}=
\frac{\mathrm{TP}\,\mathrm{TN}-\mathrm{FP}\,\mathrm{FN}}
{\sqrt{(\mathrm{TP}+\mathrm{FP})(\mathrm{TP}+\mathrm{FN})
(\mathrm{TN}+\mathrm{FP})(\mathrm{TN}+\mathrm{FN})}} .
\]
For sparse graphs, MCC is often the best single-number fidelity diagnostic
because it uses all four confusion counts without being dominated by the
large number of true negatives. Precision and recall remain essential
because they identify whether an MCC change is caused mainly by false edges
or missed true edges.

\subsection{Sign-aware recovery}

The sign-aware panel asks a stricter question: did the reconstruction assign
the correct signed edge label in \(\{-,0,+\}\)? It shows signed precision,
signed recall, signed F1, sign-aware multiclass MCC, and sign-flip rate.
Wrong-sign recovered true edges count as both signed false positives and
signed false negatives. This panel distinguishes topology recovery from Dale
sign recovery.

\subsection{Recovered source type}

The source-type panel shows how many neurons were classified by the aggregate
fit as excitatory, inhibitory, or undecided. This is a neuron-level Dale
diagnostic, not an edge-level metric. A high undecided count means the fit
may have recovered some edges but did not collect enough consistent outgoing
evidence to assign clean source type.

\subsection{False discoveries}

The false-discovery panel shows realized false-discovery proportion:
\[
  \mathrm{FDP}=\frac{\mathrm{FP}}{\max(1,\mathrm{TP}+\mathrm{FP})},
\]
and compares it to the aggregate file's approximate stability-selection
bound normalized by selected-edge count. The stability bound is a useful
conservativeness diagnostic, but it is not a formal proof of FDR control for
these overlapping bags.

\subsection{Graph density}

The density panel compares true and recovered graph density:
\[
  \rho_{\mathrm{true}} =
  \frac{\#\hbox{ true off-diagonal edges}}{N(N-1)},\qquad
  \rho_{\mathrm{reco}} =
  \frac{\#\hbox{ recovered off-diagonal edges}}{N(N-1)} .
\]
This panel quickly reveals whether a good recall score is being achieved by
making the recovered graph too dense.

\subsection{Undecided-source diagnostics}

The undecided-source panel shows the fraction of neurons classified as
undecided, the fraction of recovered edges emitted by undecided sources, and
the fraction of true edges whose source was classified as undecided. This is
important because undecided neurons are treated as an abstention category:
their nonzero \texttt{A\_prune} edges still count in edge recovery, but they
are flagged as coming from sources with ambiguous Dale evidence.

% ===============================================================
\section{Plot Group \texttt{-p b}: Raw Count Bars}
\label{sec:plot-b}
% ===============================================================

Plot group \texttt{b} shows two stacked bar charts. The first chart is titled
\textit{edge is present (counts)}. For each fit tag, the stack is ordered:
\[
  \mathrm{TP}\quad \hbox{bottom}, \qquad
  \mathrm{FN}\quad \hbox{middle}, \qquad
  \mathrm{FP}\quad \hbox{top}.
\]
True negatives are intentionally omitted. In sparse graphs, the number of
true negatives can be very large and would hide the counts that matter for
reconstruction failures.

The second chart is titled \textit{edge sign is correct (counts)} and uses
the same stacked layout for signed TP, signed FN, and signed FP. This shows
how much of the remaining error comes from missed signed edges, spurious
signed edges, or wrong signs on recovered true edges.

Raw counts matter because ratio metrics can hide scale. For example, two
fits can have similar precision but very different numbers of false edges.
The count canvas is therefore the best place to inspect the absolute error
budget.

% ===============================================================
\section{Plot Group \texttt{-p c}: Selected Confusion Matrices}
\label{sec:plot-c}
% ===============================================================

Plot group \texttt{c} displays detailed confusion matrices for selected fit
tags. The current plotting call selects \texttt{tagIdxL=[2,-1]}, meaning the
third provided tag and the last provided tag. The value \(-1\) always means
the last tag.

For each selected tag, the left matrix is the edge-label confusion matrix:
rows are true labels and columns are recovered labels in the order
\(\{-,0,+\}\). It separates:
\begin{itemize}
  \item missed inhibitory and excitatory edges: true \(-\) or \(+\) recovered
  as \(0\);
  \item spurious inhibitory and excitatory edges: true \(0\) recovered as
  \(-\) or \(+\);
  \item sign flips: true \(+\) recovered as \(-\), or true \(-\) recovered as
  \(+\).
\end{itemize}

For each selected tag, the right matrix is the neuron type confusion matrix.
Rows are true source type, excitatory or inhibitory. Columns are recovered
source type: excitatory, inhibitory, or undecided. This matrix diagnoses
Dale source classification separately from edge topology.

% ===============================================================
\section{Saved Metric File}
\label{sec:output}
% ===============================================================

The script always saves the computed metric arrays to
\texttt{edgeReco/\{outName\}.npz}. The default output name is
\texttt{\{fitNameTrunk\}\_ermHASH4}. The saved arrays include:
\begin{itemize}
  \item fit names, fit tags, input files, and truth file;
  \item duration, number of neurons, number of bags, bag fraction, FDR
  quantile, stability threshold, and \(\varepsilon\);
  \item TP, FP, FN, TN, precision, recall, F1, MCC, FDP, graph densities;
  \item signed precision, signed recall, signed F1, signed MCC, sign-flip
  count, and sign-flip rate;
  \item edge-label confusion matrices and neuron-type confusion matrices for
  every tag.
\end{itemize}

% ===============================================================
\section{Possible Future Upgrade}
\label{sec:future}
% ===============================================================

A threshold-free precision--recall study could be added later. The most
natural candidate score is not final \(|\texttt{A\_prune}|\), because
\texttt{A\_prune} has already passed per-bag FDR, stability selection, and
Dale pruning. Better future scores may come from \texttt{selection\_frequency},
\texttt{A\_mean\_selected}, null thresholds, or per-edge diagnostics saved
before the final stability cut. Such an upgrade would complement the current
absolute operating-point metrics rather than replace them.

\end{document}
